%========================================
% MatinBook - Stage 07: Boxes Test
% Test: All 10 Colored Box Types + Visual Comparison
% Purpose: Validate all tcolorbox environments and color families
%========================================

% Add parent directory to search path
\makeatletter
\def\input@path{{../}}
\makeatother

\documentclass{matinbook}

\title{آزمون جعبه‌های رنگی}
\author{تیم توسعه MatinBook}
\date{تابستان ۱۴۰۵}

\begin{document}

%----------------------------------------
% Title Page
%----------------------------------------
\maketitle

%----------------------------------------
% Chapter 1: Blue Family (Theorems)
%----------------------------------------
\chapter{خانواده آبی: قضایا و اثبات‌ها}

\section{راهنمای رنگ‌ها}

خانواده آبی برای مفاهیم اصلی و اثبات‌شده استفاده می‌شود.
شامل چهار نوع جعبه است که از تیره به روشن مرتب شده‌اند:

\begin{itemize}
    \item \colorbox{matindarkblue!80}{\textcolor{white}{\textbf{گزاره}}} --- تیره‌ترین
    \item \colorbox{matinblue}{\textcolor{white}{\textbf{قضیه}}} --- اصلی
    \item \colorbox{matinblue!70}{\textcolor{white}{\textbf{لم}}} --- متوسط
    \item \colorbox{matinblue!50}{\textcolor{white}{\textbf{نتیجه}}} --- روشن
\end{itemize}

\section{نمونه‌های خانواده آبی}

\begin{theorem}[قضیه فیثاغورث]
    در مثلث قائم‌الزاویه، مربع وتر برابر است با
    مجموع مربعات دو ضلع دیگر:
    \[
        a^{2} + b^{2} = c^{2}
    \]
    \label{thm:pythagoras}
\end{theorem}

\begin{lemma}[لم کمکی]
    این یک لم با رنگ آبی متوسط است.
    لم‌ها معمولاً برای اثبات قضایای بزرگتر
    به کار می‌روند و ارزش مستقلی ندارند.
    \label{lem:sample}
\end{lemma}

\begin{corollary}[نتیجه مستقیم]
    این یک نتیجه با رنگ آبی روشن است.
    نتایج معمولاً بلافاصله از قضایا یا لم‌ها
    نتیجه‌گیری می‌شوند.
    \label{cor:sample}
\end{corollary}

\begin{proposition}[گزاره نمونه]
    این یک گزاره با رنگ آبی تیره است.
    گزاره‌ها احکامی هستند که ارزش اثبات دارند
    اما لزوماً به اندازه قضایا مهم نیستند.
    \label{prop:sample}
\end{proposition}

\begin{proof}[اثبات نمونه]
    این یک جعبه اثبات است که از خانواده آبی
    با رنگ بسیار روشن استفاده می‌کند.
    اثبات‌ها معمولاً با علامت پایان مشخص می‌شوند. $\square$
\end{proof}

%----------------------------------------
% Chapter 2: Green Family (Definitions)
%----------------------------------------
\chapter{خانواده سبز: تعاریف و مفاهیم}

\section{راهنمای رنگ}

خانواده سبز برای تعاریف و معرفی مفاهیم جدید
استفاده می‌شود. رنگ سبز نماد رشد و یادگیری است.

\colorbox{matingreen}{\textcolor{white}{\textbf{تعریف}}} ---
رنگ اصلی خانواده سبز

\section{نمونه‌های تعاریف}

\begin{definition}[مشتق تابع]
    مشتق تابع $f$ در نقطه $x_{0}$ به صورت زیر تعریف می‌شود:
    \[
        f'(x_{0}) = \lim_{h \to 0} \frac{f(x_{0} + h) - f(x_{0})}{h}
    \]
    مشتق نشان‌دهنده نرخ تغییرات لحظه‌ای تابع است.
    \label{def:derivative}
\end{definition}

\begin{definition}[انتگرال معین]
    انتگرال معین تابع $f$ روی بازه $[a, b]$:
    \[
        \int_{a}^{b} f(x) \dd{x} = \lim_{n \to \infty}
        \sum_{i=1}^{n} f(x_{i}^{*}) \Delta x
    \]
    انتگرال سطح زیر نمودار تابع را محاسبه می‌کند.
    \label{def:integral}
\end{definition}

\begin{definition}[ماتریس]
    ماتریس $A$ با ابعاد $m \times n$ آرایه‌ای مستطیلی
    از اعداد با $m$ سطر و $n$ ستون است:
    \[
        A = [a_{ij}]_{m \times n}
    \]
    \label{def:matrix}
\end{definition}

%----------------------------------------
% Chapter 3: Orange Family (Examples)
%----------------------------------------
\chapter{خانواده نارنجی: مثال‌ها}

\section{راهنمای رنگ}

خانواده نارنجی برای مثال‌های کاربردی استفاده می‌شود.
رنگ گرم نارنجی توجه خواننده را به کاربرد مفاهیم جلب می‌کند.

\colorbox{matinorange}{\textcolor{white}{\textbf{مثال}}} ---
رنگ اصلی خانواده نارنجی

\section{نمونه‌های مثال}

\begin{example}[مشتق چندجمله‌ای]
    مشتق تابع $f(x) = x^{3} - 2x^{2} + 5x - 1$:
    \[
        f'(x) = 3x^{2} - 4x + 5
    \]
    \label{ex:derivative-poly}
\end{example}

\begin{example}[محاسبه انتگرال]
    مساحت زیر نمودار $f(x) = x^{2}$ از $x = 0$ تا $x = 2$:
    \[
        \int_{0}^{2} x^{2} \dd{x}
        = \left[ \frac{x^{3}}{3} \right]_{0}^{2}
        = \frac{8}{3} \approx 2.67
    \]
    \label{ex:integral-x2}
\end{example}

\begin{example}[ضرب ماتریس‌ها]
    \[
    \mat{1 & 2 \\ 3 & 4} \times
    \mat{5 & 6 \\ 7 & 8} =
    \mat{1\cdot5+2\cdot7 & 1\cdot6+2\cdot8 \\
         3\cdot5+4\cdot7 & 3\cdot6+4\cdot8} =
    \mat{19 & 22 \\ 43 & 50}
    \]
    \label{ex:matrix-multiply}
\end{example}

%----------------------------------------
% Chapter 4: Red Family (Remarks)
%----------------------------------------
\chapter{خانواده قرمز: نکات و هشدارها}

\section{راهنمای رنگ}

خانواده قرمز برای نکات مهم، هشدارها و توجه‌ها
استفاده می‌شود. رنگ قرمز به طور طبیعی جلب توجه می‌کند.

\colorbox{matinred}{\textcolor{white}{\textbf{نکته}}} ---
رنگ اصلی خانواده قرمز

\section{نمونه‌های نکات}

\begin{remark}[نکته مهم]
    در محاسبه مشتق، توجه داشته باشید که
    $\deriv{}{x}(x^{n}) = nx^{n-1}$ فقط برای
    $n \neq 0$ معتبر است. برای $n = 0$،
    مشتق تابع ثابت صفر است.
    \label{rem:derivative-power}
\end{remark}

\begin{remark}[هشدار]
    تقسیم بر صفر تعریف نشده است!
    هنگام ساده‌سازی عبارات کسری،
    همواره فرض کنید مخرج غیرصفر است.
    \label{rem:division-zero}
\end{remark}

\begin{remark}[نکته کاربردی]
    برای حفظ کردن فرمول مشتق توابع مثلثاتی،
    به یاد داشته باشید که مشتق سینوس، کسینوس
    و مشتق کسینوس، \textbf{منفی} سینوس است:
    \[
        \deriv{}{x}(\sin x) = \cos x, \qquad
        \deriv{}{x}(\cos x) = -\sin x
    \]
    \label{rem:trig-derivatives}
\end{remark}

%----------------------------------------
% Chapter 5: Purple Family (Exercises)
%----------------------------------------
\chapter{خانواده بنفش: تمرین‌ها}

\section{راهنمای رنگ}

خانواده بنفش برای تمرین‌ها و مسائل استفاده می‌شود.
رنگ بنفش محرک خلاقیت و تفکر است.

\colorbox{matinpurple}{\textcolor{white}{\textbf{تمرین}}} ---
رنگ اصلی خانواده بنفش

\section{نمونه‌های تمرین}

\begin{exercise}[مشتق توابع مثلثاتی]
    مشتق توابع زیر را محاسبه کنید:
    \begin{enumerate}
        \item $f(x) = \sin(2x)$
        \item $g(x) = \cos(x^{2})$
        \item $h(x) = \tan x$
    \end{enumerate}
    \label{exc:trig-derivatives}
\end{exercise}

\begin{solution}[راه حل]
    با استفاده از قاعده زنجیره‌ای:
    \begin{enumerate}
        \item $f'(x) = 2\cos(2x)$
        \item $g'(x) = -2x\sin(x^{2})$
        \item $h'(x) = \sec^{2} x = \frac{1}{\cos^{2} x}$
    \end{enumerate}
\end{solution}

\begin{exercise}[انتگرال‌گیری]
    انتگرال‌های زیر را محاسبه کنید:
    \begin{enumerate}
        \item $\int (3x^{2} - 4x + 1) \dd{x}$
        \item $\int \sin x \dd{x}$
        \item $\int e^{2x} \dd{x}$
    \end{enumerate}
    \label{exc:integrals}
\end{exercise}

\begin{solution}[راه حل]
    \begin{enumerate}
        \item $\int (3x^{2} - 4x + 1) \dd{x}
        = x^{3} - 2x^{2} + x + C$
        \item $\int \sin x \dd{x} = -\cos x + C$
        \item $\int e^{2x} \dd{x}
        = \frac{1}{2}e^{2x} + C$
    \end{enumerate}
\end{solution}

%----------------------------------------
% Chapter 6: All Boxes at a Glance
%----------------------------------------
\chapter{نمای کلی همه جعبه‌ها}

\section{مقایسه بصری}

در این بخش، هر ۱۰ نوع جعبه را بدون محتوای طولانی
در کنار هم می‌بینیم تا تفاوت رنگ‌ها مشهود باشد:

\begin{theorem}
    \textbf{قضیه (آبی)}: مهم‌ترین احکام ریاضی
\end{theorem}

\begin{lemma}
    \textbf{لم (آبی متوسط)}: احکام کمکی برای اثبات
\end{lemma}

\begin{corollary}
    \textbf{نتیجه (آبی روشن)}: نتایج مستقیم قضایا
\end{corollary}

\begin{proposition}
    \textbf{گزاره (آبی تیره)}: احکام با اهمیت کمتر
\end{proposition}

\begin{definition}
    \textbf{تعریف (سبز)}: معرفی مفاهیم جدید
\end{definition}

\begin{example}
    \textbf{مثال (نارنجی)}: کاربرد مفاهیم
\end{example}

\begin{remark}
    \textbf{نکته (قرمز)}: توضیحات و هشدارهای مهم
\end{remark}

\begin{exercise}
    \textbf{تمرین (بنفش)}: مسائل برای حل کردن
\end{exercise}

\begin{solution}
    \textbf{راه حل (سبز روشن)}: پاسخ تشریحی تمرین‌ها
\end{solution}

\begin{proof}
    \textbf{اثبات (آبی خیلی روشن)}: استدلال ریاضی
\end{proof}

%----------------------------------------
% Summary
%----------------------------------------
\chapter{خلاصه نتایج}

\section{وضعیت جعبه‌های رنگی}

\begin{table}[h]
\centering
\caption{وضعیت تمام جعبه‌های رنگی MatinBook}
\label{tab:boxes-status}
\begin{tabular}{|c|C{3cm}|c|c|}
\hline
\textbf{خانواده} & \textbf{نوع جعبه} & \textbf{رنگ} & \textbf{وضعیت} \\
\hline
\multirow{5}{*}{آبی} & قضیه (\lr{theorem}) & \colorbox{matinblue}{\textcolor{white}{آبی}} & \textcolor{matingreen}{✅} \\
& لم (\lr{lemma}) & \colorbox{matinblue!70}{\textcolor{white}{متوسط}} & \textcolor{matingreen}{✅} \\
& نتیجه (\lr{corollary}) & \colorbox{matinblue!50}{\textcolor{white}{روشن}} & \textcolor{matingreen}{✅} \\
& گزاره (\lr{proposition}) & \colorbox{matindarkblue}{\textcolor{white}{تیره}} & \textcolor{matingreen}{✅} \\
& اثبات (\lr{proof}) & \colorbox{matinblue!30}{خیلی روشن} & \textcolor{matingreen}{✅} \\
\hline
سبز & تعریف (\lr{definition}) & \colorbox{matingreen}{\textcolor{white}{سبز}} & \textcolor{matingreen}{✅} \\
\hline
نارنجی & مثال (\lr{example}) & \colorbox{matinorange}{\textcolor{white}{نارنجی}} & \textcolor{matingreen}{✅} \\
\hline
قرمز & نکته (\lr{remark}) & \colorbox{matinred}{\textcolor{white}{قرمز}} & \textcolor{matingreen}{✅} \\
\hline
بنفش & تمرین (\lr{exercise}) & \colorbox{matinpurple}{\textcolor{white}{بنفش}} & \textcolor{matingreen}{✅} \\
\hline
سبز روشن & راه حل (\lr{solution}) & \colorbox{matingreen!60}{\textcolor{white}{روشن}} & \textcolor{matingreen}{✅} \\
\hline
\end{tabular}
\end{table}

\section{نتیجه}

\textbf{تمام ۱۰ نوع جعبه رنگی \lr{MatinBook} نسخه ۱ با موفقیت آزمایش شدند.}
رنگ‌ها در ۵ خانواده (آبی، سبز، نارنجی، قرمز، بنفش)
به درستی کار می‌کنند و تمایز بصری مناسبی ایجاد می‌کنند. 🎉

\end{document}